\section{Introduction}
\label{Intro}

Logic network based quantum circuit techniques such as LUT-based Hierarchical Reversible
Synthesis (LHRS)~\cite{bib-lhrs-soeken} and \emph{caterpillar}~\cite{bib-meuli-ros}, which
uses XOR-AND-Inverter graphs (XAG) to represent the Boolean functions. XAG based methods
in particular have a range of techniques to optimize the generated quantum circuit for
T-count, which is the number of T-gates in a quantum circuit, and T-depth, which is the
largest such chain of T-gates in a quantum circuit. These metrics are important to fault
tolerant quantum computing because T-gates have an outsize implementation cost compared
to other gates~\cite{bib-bravyi-kitaev-magic-distillation}. In XAGs, T-depth can be
analogized to AND operations, which are often implemented using other three qubit
non-Clifford gates such as the Toffoli gate and the three qubit controlled-Z gate.
Therefore existing methods to optimize them for AND nodes~\cite{bib-meuli-mult} and
AND-depth\cite{bib-meuli-ros}, which obliquely optimize for T-count and T-depth
respectively.

Existing methods, however, largely use single-target gates acting on clean ancilla,
and are required to be both initialized to $\ket{0}$ and returned to the value at the end
of the calculation. Multiple Control Toffoli (MCT) gates can be decomposed with linear cost
using such ancilla. We can model such a decomposition as an XAG, one of two-input ANDs
chained together. In a similar way, we can also decompose the MCT using \emph{dirty ancilla},
which have no requirements on its initial value. Using dirty ancilla means the circuit
can use any of the other qubits that are presently being used by other subccircuits as
scratch space. This reduces the need for adding clean ancilla towards the total T-count.
In the standard formulation, dirty ancilla decompositions have higher T-count than the
clean ancilla decompositions. Some work optimizes these using a \emph{relative phase},
or a phase on the output that is input dependent, but there remains a linear separation
between the two.

In this work, we derive a way to generate XAG using dirty ancilla, inspired by the dirty
ancilla MCT decomposition. We then use relative phase to further reduce the T-count of
our decomposition, and demonstrate that it is only separated by a T-count equal to the
number of output qubits. This implies that a clean ancilla and dirty ancilla MCT
decomposition only differs by a single T-gate. We conduct several experiments on
EPFL benchmarks to demonstrate this.

This paper is structured as follows. In~\secref{Pre}, we first introduce preliminary
knowledge to understand our proposals, including XAG quantum circuit synthesis and
relative phase. We then present our proposal in~\secref{Pro}, including the dirty
ancilla XAG synthesis, and its implications with MCT decomposition. Finally, \secref{Conc}
concludes with our conclusions and future work.




%% Multiple Control Toffoli (MCT) gates are important components to quantum computing. They
%% allow non-trivial non-Clifford entanglements, which serve as building blocks to quantum
%% algorithms which demonstrate quantum advantage.

%% The canonical Toffoli gate has two control bits and a target bit. It implements the Pauli X
%% operation on its target bit when both of its input qubits are 1. A Toffoli gate is not a
%% gate that can be implemented directly; it is often decomposed into the Clifford+T basis
%% gate set. For fault tolerant computing, the T-gate dominates the implementation cost.
%% Thus, T-count, or the number of T-gates and T$^{\dagger}$-gates in a circuit, is a metric
%% that is often targeted in optimization

%% An MCT gate generalizes the Toffoli gate to multiple control bits. In general, implementing
%% the MCT gate using only the Clifford+T basis gate set without any additional bits takes
%% a T-count that is at best quadratic in the control bits~\cite{bib-barenco-elementary}.
%% However, when ancilla bits are employed, the MCT gates can be decomposed using a linear
%% number of Toffoli gates, and therefore a linear T-count~\cite{bib-barenco-elementary,bib-mike-and-ike}.
%% The methods differ whether the ancilla used are {\it clean ancilla}, which are initialized to $\ket{0}$,
%% or {\it dirty ancilla}, which have no requirements on initial values. Dirty ancilla decompositions
%% are important because they make possible implementations that simply use existing qubits
%% for ancilla, obviating the need to add qubits initialized to $\ket{0}$.

%% Decomposing MCT gates using clean ancilla typically requires a smaller T-count than decomposing
%% using dirty ancilla. In the canonical Barenco et al. formulation, the dirty ancilla decomposition
%% takes almost twice the number of Toffoli as a clean ancilla decomposition, implying a T-count
%% of the same scale. Successive results such as Maslov~\cite{bib-maslov-rtof} brought implied that
%% a T-count less than seven multiplied by the number of Toffoli gates is needed, by using
%% {\it relative phase} constructions. Other results by Abdessaieh et al~\cite{bib-abdessaied-mct}
%% close the gap between the clean ancilla and dirty ancilla decompositions by employing further
%% relative phase style constructions. However, this still has a separation in cost that is
%% linear.

%% In this work, we demonstrate that, by exploiting further relative phase opportunities, clean ancilla
%% and dirty ancilla decompositions have costs that are separated by a constant amount, that being a
%% single T-gate. This decreases the opportunity cost of using the the dirty ancilla decomposition. We
%% use it in this work to improve existing constructions that use the dirty ancilla construction,
%% including the linear cost, single dirty ancilla decomposition by Zindorf and
%% Bose~\cite{bib-zindorf-mct}. We obtain new cost calculations for such decompositions. We summarize
%% the existing state of the art of MCT decompositions in the last chapter.

%% Dctodo
%% This paper is structured as follows: \secref{Pre} goes into the necessary preliminary for this
%% work, including relative phase constructions. \secref{MCT} goes into best known methods of MCT
%% decomposition, many of them using relative phase constructions. \secref{Pro} describes how we
%% derive our proposed dirty ancilla decomposition. \secref{App} shows our application of this
%% proposal to existing methods. \secref{Conc} concludes the paper with our takeaways and possible
%% future research paths.
